% Archive record 12, line 1.
\hypertarget{part-2-12-0001-ii-a-risk-and-generalization}{%
\section{II-A --- Risk and Generalization}\label{part-2-12-0001-ii-a-risk-and-generalization}}

我们正式进入 Part II。按照前面调整过的学习方式，每次只推进一个相对独立的理论模块，优先掌握结论、成立条件及其研究意义，数学推导作为辅助。

\begin{readertable}
\begin{longtable}[]{@{}
  >{\raggedright\arraybackslash}p{(\columnwidth - 4\tabcolsep) * \real{0.16}}
  >{\raggedright\arraybackslash}p{(\columnwidth - 4\tabcolsep) * \real{0.62}}
  >{\raggedright\arraybackslash}p{(\columnwidth - 4\tabcolsep) * \real{0.18}}@{}}
\toprule
Section & Topics & Progress \\
\midrule
\endhead
II-A & Population Risk, Empirical Risk, Generalization & 当前 \\
II-B & Concentration Inequalities, Uniform Convergence & 待学习 \\
II-C & Model Selection Bias, Multiple Testing & 待学习 \\
II-D & Adaptive Overfitting, Holdout Reuse & 待学习 \\
II-E & Data Leakage, Distribution Shift & 待学习 \\
II-F & Time-Series Cross-Validation, Backtesting & 待学习 \\
\bottomrule
\end{longtable}
\end{readertable}

\hypertarget{part-2-12-0001-research-motivation}{%
\subsection{1. Research Motivation}\label{part-2-12-0001-research-motivation}}

Statistical Learning Theory 最根本的问题是：

\textbf{我们如何利用有限的观测数据，学习一个能够在未知数据上有效工作的模型？}

这涉及两个不同的优化目标：

\begin{itemize}
\tightlist
\item
  在已经观察到的数据上取得良好表现；
\item
  在尚未观察到的数据上取得良好表现。
\end{itemize}

前者可以直接计算，后者原则上只能估计。

这两者之间的差异，就是 Statistical Learning Theory 的重要研究对象，也是 X Program \#161 所研究的 Generalization Decay 的数学基础。

\hypertarget{part-2-12-0001-mathematical-framework}{%
\subsection{2. Mathematical Framework}\label{part-2-12-0001-mathematical-framework}}

\hypertarget{part-2-12-0001-data-generating-distribution}{%
\subsubsection{2.1 Data Generating Distribution}\label{part-2-12-0001-data-generating-distribution}}

设随机变量：
\[Z=(X,Y)\sim\mathcal D,\]

其中：

\begin{itemize}
\tightlist
\item
  \(X\in\mathcal X\)：Input / Feature。
\item
  \(Y\in\mathcal Y\)：Target / Label。
\item
  \(\mathcal D\)：未知的联合数据生成分布。
\item
  \(h:\mathcal X\to\mathcal Y\)：Hypothesis，即一个预测模型。
\end{itemize}

进一步定义：
\[\mathcal H=\{h:h\text{ is an admissible model}\},\]

称为 \textbf{Hypothesis Space（假设空间）}。

例如，不同的 Hypothesis Space 可以是：

\begin{readertable}
\begin{longtable}[]{@{}
  >{\raggedright\arraybackslash}p{(\columnwidth - 2\tabcolsep) * \real{0.42}}
  >{\raggedright\arraybackslash}p{(\columnwidth - 2\tabcolsep) * \real{0.54}}@{}}
\toprule
Model Family & Hypothesis Space \\
\midrule
\endhead
Linear Regression & 所有允许的线性函数 \\
Polynomial Regression & 所有满足阶数约束的多项式函数 \\
Decision Tree & 所有满足深度等约束的决策树 \\
Neural Network & 所有满足结构约束的神经网络 \\
AI Agent & 所有允许的 Agent 程序或工作流程 \\
\bottomrule
\end{longtable}
\end{readertable}

最后一个是对传统学习框架的推广。在 RSI 场景中，Agent 本身也可以被视为一个需要评价和选择的 Hypothesis。

\hypertarget{part-2-12-0001-loss-function}{%
\subsubsection{2.2 Loss Function}\label{part-2-12-0001-loss-function}}

定义损失函数：
\[\ell(h(X),Y).\]

它用于衡量模型预测结果与真实结果之间的差异。

例如：

\begin{readertable}
\begin{longtable}[]{@{}
  >{\raggedright\arraybackslash}p{(\columnwidth - 2\tabcolsep) * \real{0.43}}
  >{\raggedright\arraybackslash}p{(\columnwidth - 2\tabcolsep) * \real{0.53}}@{}}
\toprule
Loss Function & Mathematical Expression \\
\midrule
\endhead
Squared Error & \((h(X)-Y)^2\) \\
Absolute Error & \(\lvert h(X)-Y\rvert\) \\
Zero-One Loss & \(\mathbf 1_{\{h(X)\neq Y\}}\) \\
Binary Cross-Entropy & \(-Y\log p-(1-Y)\log(1-p)\) \\
\bottomrule
\end{longtable}
\end{readertable}

其中 \(p\) 是模型预测 \(Y=1\) 的概率。

从这里开始，我们暂时统一采用 \textbf{Risk（越低越好）} 的记号，而不是 Part I 中的 Performance Score（越高越好）。

对于能够表示为样本平均值的性能指标，两者可以通过令 Loss 等于负 Score 进行转换。但 Sharpe Ratio 等非线性总体指标需要单独处理。

\hypertarget{part-2-12-0001-population-risk-and-empirical-risk}{%
\subsection{3. Population Risk and Empirical Risk}\label{part-2-12-0001-population-risk-and-empirical-risk}}

这是整个 Part II 最基础的一组定义。

\hypertarget{part-2-12-0001-population-riskux603bux4f53ux98ceux9669}{%
\subsubsection{3.1 Population Risk（总体风险）}\label{part-2-12-0001-population-riskux603bux4f53ux98ceux9669}}

定义：
\[\boxed{
R(h)=
\mathbb E_{(X,Y)\sim\mathcal D}
[\ell(h(X),Y)]
}\]

如果展开为积分，则有：
\[R(h)=
\int_{\mathcal X\times\mathcal Y}
\ell(h(x),y)\,\mathrm d\mathcal D(x,y).\]

Population Risk 描述的是模型在目标数据分布上的真实期望损失。

因此，我们真正希望解决的优化问题是：
\[\boxed{
h^*=\arg\min_{h\in\mathcal H}R(h)
}\]

其中假设最优解存在。若不存在，则应使用 infimum 描述最优风险。

这里的根本困难是：我们通常并不知道 \(\mathcal D\)。

所以，Population Risk 虽然是我们真正关心的量，却通常无法直接计算。

\hypertarget{part-2-12-0001-empirical-riskux7ecfux9a8cux98ceux9669}{%
\subsubsection{3.2 Empirical Risk（经验风险）}\label{part-2-12-0001-empirical-riskux7ecfux9a8cux98ceux9669}}

假设我们观察到 \(n\) 个独立同分布的样本：
\[D_n=\{(X_i,Y_i)\}_{i=1}^{n}.\]

定义：
\[\boxed{
\widehat R_n(h)
=
\frac1n
\sum_{i=1}^{n}
\ell(h(X_i),Y_i)
}\]

这就是 Empirical Risk。

其本质是利用 Sample Mean 估计 Population Mean：
\[\underbrace{
\mathbb E[\ell(h(X),Y)]
}_{\text{Population Risk}}
\quad\longleftarrow\quad
\underbrace{
\frac1n\sum_{i=1}^{n}\ell(h(X_i),Y_i)
}_{\text{Empirical Risk}}.\]

\bookgraphic{0.531}{part-2-12-0001-001-01}{}

\textbf{Population Risk} \(R(h)\)

真正希望最小化的量

\textbf{Empirical Risk} \(\widehat R_n(h)\)

实际可以计算的量

\hypertarget{part-2-12-0001-empirical-risk-minimization-erm}{%
\subsubsection{3.3 Empirical Risk Minimization (ERM)}\label{part-2-12-0001-empirical-risk-minimization-erm}}

由于无法直接最小化 Population Risk，我们转而求解：
\[\boxed{
\hat h_n
=
\arg\min_{h\in\mathcal H}
\widehat R_n(h)
}\]

这就是 \textbf{Empirical Risk Minimization（经验风险最小化）}。

实际上，我们之前学习的很多统计方法都是 ERM 的具体实例。

例如 Ordinary Least Squares：
\[\hat{\boldsymbol\beta}
=
\arg\min_{\boldsymbol\beta}
\frac1n
\sum_{i=1}^{n}
(y_i-\mathbf x_i^\top\boldsymbol\beta)^2.\]

它最小化的就是 Squared Loss 对应的 Empirical Risk。

而 Statistical Learning Theory 关心的是：
\[\boxed{
\hat h_n\text{ minimizes }\widehat R_n(h)
\quad\Longrightarrow\quad
\text{How good is }R(\hat h_n)?
}\]

注意，这里不是一个已经成立的数学蕴含，而是我们试图回答的研究问题。

\hypertarget{part-2-12-0001-generalization-and-generalization-gap}{%
\subsection{4. Generalization and Generalization Gap}\label{part-2-12-0001-generalization-and-generalization-gap}}

\hypertarget{part-2-12-0001-generalization-gap}{%
\subsubsection{4.1 Generalization Gap}\label{part-2-12-0001-generalization-gap}}

由于我们现在采用 Risk 记号，定义：
\[\boxed{
G_n(h)=R(h)-\widehat R_n(h)
}\]

\begin{readertable}
\begin{longtable}[]{@{}
  >{\raggedright\arraybackslash}p{(\columnwidth - 2\tabcolsep) * \real{0.31}}
  >{\raggedright\arraybackslash}p{(\columnwidth - 2\tabcolsep) * \real{0.65}}@{}}
\toprule
Gap & Interpretation \\
\midrule
\endhead
\(G_n(h)>0\) & 经验损失低估了真实损失 \\
\(G_n(h)=0\) & 经验损失与真实损失一致 \\
\(G_n(h)<0\) & 经验损失高估了真实损失 \\
\bottomrule
\end{longtable}
\end{readertable}

正向的 Generalization Gap 通常对应我们关心的过度乐观现象。

但必须注意：

\textbf{Generalization Gap 小，并不意味着模型性能好。}

例如：

假设以下数值为模型的 MSE，越低越好。

\textbf{Model A --- Underfitted}

Gap = 0.02

Empirical

0.80

Population

0.82

\textbf{Model B --- Well-generalized}

Gap = 0.04

Empirical

0.20

Population

0.24

\textbf{Model C --- Overfitted}

Gap = 0.60

Empirical

0.05

Population

0.65

Model A 的 Generalization Gap 最小，但其真实性能最差。

因此，我们真正的优化目标仍然是：
\[\min R(h),\]

而不是：
\[\min G_n(h).\]

\textbf{Research Implication：} X Program \#161 中提出的 Mitigation Mechanism 不能仅仅以缩小 IS/OOS Gap 作为成功标准。

如果一个机制通过降低所有版本的 IS Performance 来缩小 Gap，却没有改善 OOS Performance，那么它并没有真正解决我们关心的性能问题。

\hypertarget{part-2-12-0001-generalization-ux4e0e-model-complexity}{%
\subsubsection{4.2 Generalization 与 Model Complexity}\label{part-2-12-0001-generalization-ux4e0e-model-complexity}}

传统 Statistical Learning 中一个典型的现象是：随着 Model Complexity 增加，Training Risk 通常下降，而 Population Risk 可能先下降后上升。

\textbf{Classical Underfitting--Overfitting Curve}

\bookgraphic{0.487}{part-2-12-0001-003-01}{}

Empirical Risk

Population Risk

经典理论示意，并非所有模型都必然呈现 U-shaped Risk Curve。

你应该能够立即发现：这与 X Program \#161 的 Better, Then Worse 现象非常相似。

但存在一个重要区别。

\begin{readertable}
\begin{longtable}[]{@{}
  >{\raggedright\arraybackslash}p{(\columnwidth - 2\tabcolsep) * \real{0.41}}
  >{\raggedright\arraybackslash}p{(\columnwidth - 2\tabcolsep) * \real{0.55}}@{}}
\toprule
Classical Overfitting & RSI Performance Decay \\
\midrule
\endhead
横轴是 Model Complexity & 横轴是 Self-Improvement Iteration \\
主要关注模型拟合 & 同时关注模型拟合与自适应选择 \\
优化空间通常事先定义 & 优化空间可能不断变化 \\
Training Data 可能被过拟合 & 整个 Evaluation Environment 可能被过拟合 \\
\bottomrule
\end{longtable}
\end{readertable}

在 RSI 中，Iteration 与 Model Complexity 不一定同步增长。

例如，一个 Agent 后期可能删减代码、降低模型参数量，却通过更有针对性的搜索策略加剧 Adaptive Overfitting。

因此，不能简单地认为：
\[\text{More Iterations}
\equiv
\text{Higher Model Complexity}.\]

这两者之间的关系本身就可以成为研究问题。

\hypertarget{part-2-12-0001-pointwise-convergence-vs.-uniform-convergence}{%
\subsection{5. Pointwise Convergence vs. Uniform Convergence}\label{part-2-12-0001-pointwise-convergence-vs.-uniform-convergence}}

这一节是本模块最重要的理论内容，也是下一模块的入口。

\hypertarget{part-2-12-0001-pointwise-convergence}{%
\subsubsection{5.1 Pointwise Convergence}\label{part-2-12-0001-pointwise-convergence}}

利用我们已经学过的 Law of Large Numbers。

对于任意一个事先固定的模型 \(h\)，若损失可积，且样本独立同分布，则根据 SLLN：
\[\boxed{
\widehat R_n(h)
\xrightarrow{\mathrm{a.s.}}
R(h)
}\]

也就是说，对于一个固定模型，随着样本量增加，Empirical Risk 几乎必然收敛于 Population Risk。

我们称之为 \textbf{Pointwise Convergence（逐点收敛）}。

这个结论本身并不复杂。

但问题在于，我们真正关心的不是任意一个固定模型，而是通过数据选择出来的模型：
\[\hat h_n
=
\arg\min_{h\in\mathcal H}
\widehat R_n(h).\]

这里的 \(\hat h_n\) 随数据和样本量变化。

因此，我们不能直接应用固定模型的 LLN，断言：
\[\widehat R_n(\hat h_n)
-
R(\hat h_n)
\xrightarrow{\mathrm{a.s.}}0.\]

\textbf{对于一般的无限假设空间，这个推论并不自动成立。}

\hypertarget{part-2-12-0001-uniform-convergence}{%
\subsubsection{5.2 Uniform Convergence}\label{part-2-12-0001-uniform-convergence}}

为了控制经过选择的模型，我们希望建立一个更强的结论：
\[\boxed{
\sup_{h\in\mathcal H}
|\widehat R_n(h)-R(h)|
\xrightarrow{\mathrm{a.s.}}0
}\]

这就是一种 Uniform Convergence（一致收敛）条件。

\textbf{Pointwise Convergence}

\bookgraphic{0.341}{part-2-12-0001-004-01}{}

\begin{verbatim}
  对每一个固定模型分别控制误差，但不同模型的收敛速度可能不同。
\end{verbatim}

\textbf{Uniform Convergence}

\bookgraphic{0.341}{part-2-12-0001-004-02}{}

\begin{verbatim}
  同时控制整个 Hypothesis Space 中最大的评估误差。
\end{verbatim}

两者的区别可以直接从量词顺序理解。

Pointwise Convergence：
\[\forall h\in\mathcal H,\quad
\forall\epsilon>0,\quad
\exists N(h,\epsilon).\]

Uniform Convergence：
\[\forall\epsilon>0,\quad
\exists N(\epsilon),\quad
\forall h\in\mathcal H.\]

这里省略了随机样本路径及后续 \(n\geq N\) 的条件，以突出量词顺序的区别。

在 Uniform Convergence 中，控制误差所需要的样本量不再针对某一个具体模型，而是需要同时适用于整个模型集合。

\hypertarget{part-2-12-0001-ux4e3aux4ec0ux4e48ux8fd9ux5bf9ux4e8e-rsi-ux7279ux522bux91cdux8981}{%
\paragraph{为什么这对于 RSI 特别重要？}\label{part-2-12-0001-ux4e3aux4ec0ux4e48ux8fd9ux5bf9ux4e8e-rsi-ux7279ux522bux91cdux8981}}

因为 RSI 不是在评价一个事先固定的 Agent，而是在大量候选版本中不断寻找表现最好的版本。

我们真正需要控制的是：
\[\sup_{A\in\mathcal A}
|\widehat R_n(A)-R(A)|.\]

而不是：
\[|\widehat R_n(A_0)-R(A_0)|.\]

对于固定有限的 Hypothesis Space，可以利用 LLN 获得一致收敛；但对于复杂甚至开放式的 Agent Space，问题就远没有那么简单。

下一模块我们将正式研究如何对这个 supremum 建立概率上界。

\hypertarget{part-2-12-0001-excess-risk-ux5982ux4f55ux8bc4ux4ef7ux6a21ux578bux9009ux62e9ux7684ux8d28ux91cf}{%
\subsection{6. Excess Risk --- 如何评价模型选择的质量？}\label{part-2-12-0001-excess-risk-ux5982ux4f55ux8bc4ux4ef7ux6a21ux578bux9009ux62e9ux7684ux8d28ux91cf}}

Generalization Gap 评价的是经验风险与总体风险之间的差异。

但另一个同样重要的问题是：

\textbf{我们选择出来的模型，与假设空间中真正最好的模型相比，究竟差多少？}

定义 Excess Risk：
\[\boxed{
\mathcal E(\hat h_n)
=
R(\hat h_n)-R(h^*)
}\]

其中：
\[h^*=\arg\min_{h\in\mathcal H}R(h),\]

\[\hat h_n=\arg\min_{h\in\mathcal H}\widehat R_n(h).\]

\hypertarget{part-2-12-0001-excess-risk-bound}{%
\subsubsection{6.1 Excess Risk Bound}\label{part-2-12-0001-excess-risk-bound}}

这是一个非常经典且重要的结论：
\[\boxed{
R(\hat h_n)-R(h^*)
\leq
2\sup_{h\in\mathcal H}
|\widehat R_n(h)-R(h)|
}\]

也就是说，只要能够同时控制整个假设空间的 Generalization Gap，就能够控制 ERM 所选择模型的实际性能损失。

\hypertarget{part-2-12-0001-proof}{%
\paragraph{Proof}\label{part-2-12-0001-proof}}

这一段可以暂时跳过，之后通过索引回来看。

\[\begin{aligned}
R(\hat h_n)-R(h^*)
={}&
\underbrace{
R(\hat h_n)-\widehat R_n(\hat h_n)
}_{\text{Generalization Error}}
\\
&+
\underbrace{
\widehat R_n(\hat h_n)-\widehat R_n(h^*)
}_{\text{Empirical Optimization}}
\\
&+
\underbrace{
\widehat R_n(h^*)-R(h^*)
}_{\text{Estimation Error}}.
\end{aligned}\]

由于 \(\hat h_n\) 是 Empirical Risk Minimizer：
\[\widehat R_n(\hat h_n)
\leq
\widehat R_n(h^*).\]

因此中间项非正，从而：
\[\begin{aligned}
R(\hat h_n)-R(h^*)
&\leq
|R(\hat h_n)-\widehat R_n(\hat h_n)|\\
&\quad+
|\widehat R_n(h^*)-R(h^*)|\\
&\leq
2\sup_{h\in\mathcal H}
|\widehat R_n(h)-R(h)|.
\end{aligned}\]

\[\hfill\blacksquare\]

\hypertarget{part-2-12-0001-extension-approximate-optimization}{%
\subsubsection{6.2 Extension --- Approximate Optimization}\label{part-2-12-0001-extension-approximate-optimization}}

实际上，AI Agent 通常不可能找到整个程序空间中 Empirical Risk 最低的 Agent。

假设它只能找到一个近似最优解 \(\tilde h_n\)，满足：
\[\widehat R_n(\tilde h_n)
\leq
\widehat R_n(h^*)+\delta,\]

其中 \(\delta\geq0\) 表示相对于 \(h^*\) 的经验风险优化误差上界。

通过同样的证明，可以得到：
\[\boxed{
R(\tilde h_n)-R(h^*)
\leq
2\sup_{h\in\mathcal H}
|\widehat R_n(h)-R(h)|+\delta
}\]

这个不等式尤其值得我们关注，因为它建立了 \textbf{Optimization 与 Generalization 之间的关系}。

\textbf{Excess Risk Upper Bound}

\begin{verbatim}
$R(\tilde h_n)-R(h^*)$
\end{verbatim}

↓

\textbf{Statistical Error} \(2\sup_{h\in\mathcal H}|G_n(h)|\)

有限数据造成的评估不确定性

\textbf{Optimization Error} \(\delta\)

搜索算法未找到足够好的解

这给出了我们理解 RSI 的另一种角度。

随着 Self-Improvement 进行，Agent 可能越来越善于搜索，从而降低 Optimization Error。

但是，如果这一过程同时扩大了 Hypothesis Space，或者增加了对同一评估数据的自适应利用，Statistical Error 就可能变得更难控制。

\textbf{Optimization--Generalization Trade-off}

假设性的研究机制示意，而非由上述不等式直接推出的性能曲线。

\bookgraphic{0.536}{part-2-12-0001-006-01}{}

Optimization error

Statistical error

Combined bound

不过，需要严格区分：

\begin{itemize}
\tightlist
\item
  对于一个固定的 \(\mathcal H\) 和固定的数据集，上述 supremum 本身不会因为迭代次数增加而改变。
\item
  如果有效搜索空间扩大，或者研究对象变成自适应产生的随机模型集合，我们才需要进一步分析统计误差如何变化。
\item
  Excess Risk 的上界增加，也不等于实际 Excess Risk 必然增加。
\end{itemize}

因此，Optimization--Generalization Trade-off 是一个很有价值的研究解释框架，但还不是关于 RSI Performance Decay 的完整定理。

\hypertarget{part-2-12-0001-in-sample-and-out-of-sample-evaluation}{%
\subsection{7. In-Sample and Out-of-Sample Evaluation}\label{part-2-12-0001-in-sample-and-out-of-sample-evaluation}}

\hypertarget{part-2-12-0001-ux6570ux636eux96c6ux7684ux529fux80fdux5212ux5206}{%
\subsubsection{7.1 数据集的功能划分}\label{part-2-12-0001-ux6570ux636eux96c6ux7684ux529fux80fdux5212ux5206}}

在传统机器学习中，我们通常使用三个数据集。

\textbf{Training}

\begin{verbatim}
  60%
\end{verbatim}

\textbf{Validation}

\begin{verbatim}
  20%
\end{verbatim}

\textbf{Test}

\begin{verbatim}
  20%
\end{verbatim}

\begin{readertable}
\begin{longtable}[]{@{}
  >{\raggedright\arraybackslash}p{(\columnwidth - 2\tabcolsep) * \real{0.26}}
  >{\raggedright\arraybackslash}p{(\columnwidth - 2\tabcolsep) * \real{0.70}}@{}}
\toprule
\textbf{Training} & 用于参数估计及模型拟合。 \\
\midrule
\endhead
\textbf{Validation} & 用于模型选择、超参数调整和 Agent Revision。 \\
\textbf{Test} & 用于独立评估泛化性能。 \\
\bottomrule
\end{longtable}
\end{readertable}

60/20/20 仅为划分示例，不是金融时间序列的推荐固定比例。

在 RSI 项目中，最重要的是理解 \textbf{In-Sample 的含义取决于我们正在优化什么}。

考虑一个 Agent 使用 Training Set 拟合模型，并使用 Validation Set 来决定是否接受新的 Revision。

那么：

\begin{readertable}
\begin{longtable}[]{@{}
  >{\raggedright\arraybackslash}p{(\columnwidth - 4\tabcolsep) * \real{0.22}}
  >{\raggedright\arraybackslash}p{(\columnwidth - 4\tabcolsep) * \real{0.34}}
  >{\raggedright\arraybackslash}p{(\columnwidth - 4\tabcolsep) * \real{0.40}}@{}}
\toprule
Dataset & 对预测模型而言 & 对 RSI Agent 而言 \\
\midrule
\endhead
Training & In-Sample & Development Data \\
Validation & 未用于直接参数拟合 & In-Sample Evaluation \\
Test & Out-of-Sample & Out-of-Sample Evaluation \\
\bottomrule
\end{longtable}
\end{readertable}

换句话说，某个数据集对于内层模型而言可能是 Out-of-Sample，但对于外层 Agent Optimization 而言却是 In-Sample。

\textbf{In-Sample / Out-of-Sample 应当根据数据是否参与当前层级的优化与选择来判断，而不只是根据文件名称判断。}

\hypertarget{part-2-12-0001-rsi-ux4e2dux7684-oos-evaluation}{%
\subsubsection{7.2 RSI 中的 OOS Evaluation}\label{part-2-12-0001-rsi-ux4e2dux7684-oos-evaluation}}

\bookgraphic{0.477}{part-2-12-0001-008-01}{保存所有 Agent 版本，使用独立的 OOS Evaluator 进行外部审计，但不让其结果参与 Agent 的优化过程。}

这正好对应项目要求：

\begin{quote}
records the in-sample score of every version and additionally scores each one out of sample.
\end{quote}

但这里存在一个潜在的 Methodological Problem。

假设研究人员通过 OOS Evaluation 发现 \(A_5\) 表现最好，然后设计了一种机制，要求所有 Agent 在第 5 轮停止。

那么这个 OOS Dataset 实际上也参与了 Mitigation Mechanism 的设计。

因此，它不再适合作为该机制的最终独立测试集。

一种候选实验协议是：

\begin{readertable}
\begin{longtable}[]{@{}
  >{\raggedright\arraybackslash}p{(\columnwidth - 2\tabcolsep) * \real{0.42}}
  >{\raggedright\arraybackslash}p{(\columnwidth - 2\tabcolsep) * \real{0.54}}@{}}
\toprule
Dataset & Purpose \\
\midrule
\endhead
\(D_{\mathrm{train}}\) & 训练金融模型 \\
\(D_{\mathrm{IS}}\) & 驱动 RSI \\
\(D_{\mathrm{audit}}\) & 分析退化轨迹、设计缓解机制 \\
\(D_{\mathrm{test}}\) & 独立检验最终缓解机制 \\
\bottomrule
\end{longtable}
\end{readertable}

我认为这是本节与实际科研联系最紧密的一个观点。

\textbf{不仅 Agent 会过拟合，研究 Agent 的研究人员也可能过拟合。}

当然，这种四层划分是我们提出的候选实验设计，并不是已知的课题组正式协议。有限金融数据如何在各评估阶段之间分配，也需要进一步研究。

\hypertarget{part-2-12-0001-ii-a-conclusion-index}{%
\subsection{8. II-A --- Conclusion Index}\label{part-2-12-0001-ii-a-conclusion-index}}

按照我们刚刚约定的学习方式，下面是本模块需要保留的数学结论。证明与解释可以通过索引返回相应章节。

\begin{readertable}
\begin{longtable}[]{@{}
  >{\raggedright\arraybackslash}p{(\columnwidth - 4\tabcolsep) * \real{0.12}}
  >{\raggedright\arraybackslash}p{(\columnwidth - 4\tabcolsep) * \real{0.69}}
  >{\raggedright\arraybackslash}p{(\columnwidth - 4\tabcolsep) * \real{0.15}}@{}}
\toprule
ID & Key Result & Index \\
\midrule
\endhead
A1 & Population Risk 与 Empirical Risk & §3.1--2 \\
A2 & Empirical Risk Minimization & §3.3 \\
A3 & Generalization Gap 小不等于真实性能好 & §4.1 \\
A4 & RSI Iteration 不等于 Model Complexity & §4.2 \\
A5 & Fixed-model LLN 不自动保证 Selected-model Generalization & §5.1 \\
A6 & Uniform Convergence 可以同时控制所有候选模型 & §5.2 \\
A7 & Excess Risk Bound & §6.1 \\
A8 & Optimization Error 与 Statistical Error & §6.2 \\
A9 & IS/OOS 具有 Optimization-level Dependence & §7.1 \\
A10 & OOS 数据一旦参与机制设计，就需要新的最终测试集 & §7.2 \\
\bottomrule
\end{longtable}
\end{readertable}

\hypertarget{part-2-12-0001-core-formula-sheet}{%
\subsubsection{8.1 Core Formula Sheet}\label{part-2-12-0001-core-formula-sheet}}

\begin{readertable}
\begin{longtable}[]{@{}
  >{\raggedright\arraybackslash}p{(\columnwidth - 2\tabcolsep) * \real{0.32}}
  >{\raggedright\arraybackslash}p{(\columnwidth - 2\tabcolsep) * \real{0.64}}@{}}
\toprule
Concept & Formula \\
\midrule
\endhead
Population Risk & \(R(h)=\mathbb E[\ell(h(X),Y)]\) \\
Empirical Risk & \(\widehat R_n(h)=\frac1n\sum_i\ell(h(X_i),Y_i)\) \\
ERM & \(\hat h_n=\arg\min_h\widehat R_n(h)\) \\
Generalization Gap & \(G_n(h)=R(h)-\widehat R_n(h)\) \\
Excess Risk & \(\mathcal E(\hat h_n)=R(\hat h_n)-R(h^*)\) \\
\bottomrule
\end{longtable}
\end{readertable}

最重要的三个数学结论：

\textbf{Result 1 --- Pointwise LLN}

\[\widehat R_n(h)
\xrightarrow{\mathrm{a.s.}}R(h)\]

\textbf{Result 2 --- Uniform Convergence Condition}

\[\sup_{h\in\mathcal H}
|\widehat R_n(h)-R(h)|
\xrightarrow{\mathrm{a.s.}}0\]

\textbf{Result 3 --- Approximate ERM Excess Risk Bound}

\[R(\tilde h_n)-R(h^*)
\leq
2\sup_{h\in\mathcal H}
|\widehat R_n(h)-R(h)|+\delta.\]

\hypertarget{part-2-12-0001-knowledge-checklist}{%
\subsubsection{8.2 Knowledge Checklist}\label{part-2-12-0001-knowledge-checklist}}

\textbf{II-A --- Knowledge Review}

\begin{itemize}
\item
  理解 Population Risk 与 Empirical Risk 的区别
\item
  理解 ERM 为什么不一定能够获得真正的最优模型
\item
  区分 Generalization Gap 与 Population Risk
\item
  理解 Pointwise Convergence 与 Uniform Convergence 的区别
\item
  理解 Excess Risk Bound 的数学意义
\item
  理解 Optimization Error 与 Statistical Error 的区别
\item
  理解 IS/OOS 在不同优化层级上的相对性
\item
  理解 OOS Evaluation 本身也可能受到研究人员的过拟合
\end{itemize}

\hypertarget{part-2-12-0001-next-ii-b-concentration-inequalities-and-uniform-convergence}{%
\subsection{Next --- II-B: Concentration Inequalities and Uniform Convergence}\label{part-2-12-0001-next-ii-b-concentration-inequalities-and-uniform-convergence}}

目前，我们已经知道 Uniform Convergence 为什么重要，但是还没有解决一个问题：
\[\sup_{h\in\mathcal H}
|\widehat R_n(h)-R(h)|\]

究竟有多大？

下一模块将利用你已经熟悉的 Probability Theory，沿着如下路线展开：

\textbf{Law of Large Numbers}

误差最终会不会消失？

↓ \textbf{Concentration Inequalities}

有限样本下，误差超过阈值的概率有多大？

↓ \textbf{Hoeffding's Inequality}

控制单个固定模型的评估误差

↓ \textbf{Union Bound}

将误差控制扩展到多个候选模型

↓ \textbf{Uniform Generalization Bound}

\[2|\mathcal H|e^{-2n\epsilon^2}\]

建立样本量、搜索空间与泛化误差之间的数学关系

II-B 将真正解释为什么 Search Budget 和 Evaluation Sample Size 对 RSI 的可靠性如此重要，并为之后的 Adaptive Overfitting 理论建立基准。
